% ---------------------------------------------------------------------------
% Experiment History -- subsection 7: Detector Characterisation and Errata.
% Sources: README A.6.  Jobs: 101860 (Phase 0), 101866 (Phase 0.5),
% 102115 (recheck).
%
% Rewritten 2026-09-11: cut from ~700 words to ~330, to the register of the
% user-written parts 2 and 4.  Table tab:inband carries the control, so the
% prose only has to say what it means.
%
% Carries \label{subsec:exphist:characterisation}, forward-referenced by
% parts 5 and 6.  IF THIS PART IS DROPPED OR RENAMED, BOTH DANGLE.
%
% Both errata are stated with their mechanism and neither is apologised for.
% Keep the asymmetry visible: the STFT fix left its finding standing, the
% energy detector destroyed one.
%
% Figure: artifacts/frontier/inband_vs_outofband.png, four panels, as-is.
% ---------------------------------------------------------------------------

\subsection{Detector Characterisation and Errata}
\label{subsec:exphist:characterisation}
Every attack so far had run against a detector whose only published number was its accuracy on its own
training set, which says what it separates and not what it resists.
So instead of learning another attacker against it, we enumerated it: a grid over active subcarriers,
power and temporal structure through the frozen detector, recording $\Pr[\text{flag}]$ and BER per cell.
The decisive cell is a control that fixes power and damage and varies only whether the interference
touches the twelve guard and DC bins the waveform leaves empty (Table~\ref{tab:inband}).

\begin{table}[!t]
\renewcommand{\arraystretch}{1.3}
\caption{The frozen detector at fixed power and fixed damage. Only guard and DC occupancy differs
between the first two rows.}
\label{tab:inband}
\centering
\begin{tabular}{l c c c}
\hline\hline
\textbf{Interference} & \textbf{Power} & $\Pr[\text{flag}]$ & \textbf{BER} \\
\hline
in-band, 52 subcarriers  & 8    & 0.011  & 0.360 \\
out-of-band, all 64 bins & 8    & 0.999  & 0.363 \\
out-of-band, all 64 bins & 0.03 & 0.9999 & 0.000 \\
\hline\hline
\end{tabular}
\end{table}

It fires on any energy outside the occupied band, including energy far too weak to break a single bit,
and is close to blind to in-band interference that breaks a third of them.
What reproduced the published accuracy is a spectral mask monitor, not a jamming detector, and the two
only coincide because every jammer in its training set violates the mask.
That also accounts for both walls the previous step hit: the broadband probe flagged at power $0.01$
was injecting into guard and DC bins, and top-$K$ ran over all 64 bins, so activating more subcarriers
meant selecting null ones.
Our reading of that threshold -- a fundamental tension between effectiveness and stealth, BER capped
near $0.04$ -- does not survive it.

\begin{figure*}[!t]
\centering
% \includegraphics[width=\textwidth]{fig/inband_vs_outofband.pdf}
\caption{Detection probability and BER against injected power, for interference confined to the
occupied subcarriers and for interference that also occupies the guard and DC bins.}
\label{fig:inband}
\end{figure*}

Retraining with in-band interference added closes the blind spot, moving the first row from $0.011$ to
$0.995$, and charges for it on every clean frame: accuracy $99.79 \to 90.5\,\%$, false alarms
$0.00 \to 3.81\,\%$, fresh in-band jammers caught only $68\,\%$ of the time, and the omniscient
counter-signal, whose received distribution is exactly a clean one, now flagged at $0.15$.
The CNN cannot separate spectrally compliant in-band interference from a clean frame without spending
false alarms, because at the resolution it sees them the two differ by a scale factor.
That is a limit of the representation rather than a gap in the training set, and it is the first
measurement of what an adaptation costs the defender.

Two corrections belong here, and they did not cost the same.
The transform had been taken over the in-phase component alone instead of the complex two-sided
transform the reference uses; fixing it and retraining leaves the replication at $99.8\,\%$ accuracy and
zero false alarms, and leaves the control intact at $0.086$ in-band against $1.000$ out-of-band.
The second removes a result of ours: the sweep scored the CNN alone, and the receiver is not obliged to
use only the CNN.
Adding an energy detector at $1\,\%$ false alarms and flagging if either fires drops the best stealthy
BER from $0.419$ to $0.005$, and to nothing at $\Pr[\text{flag}] \leq 0.05$.
In-band interference that flips bits has to raise the received power, and a power meter is free.

What keeps the question open is that this channel is lossless, so clean power is a sharp constant and
anything added to it crosses the threshold -- which is a fact about the channel, not about jamming.
